% =====================================================================
\part{A matemática, método a método}
% =====================================================================

Esta parte é o coração do caderno. Cada capítulo segue a mesma estrutura: a
pergunta que o método responde, a fórmula, a tradução em português corrente, a
fronteira do que ele autoriza afirmar, e a formulação pronta para a banca.

Um princípio atravessa tudo o que vem a seguir. \textbf{Nenhum destes métodos
estabelece causalidade.} O que muda de um para outro é o que cada um consegue
descartar. Saber exatamente o que cada instrumento descarta --- e o que ele
deixa de pé --- é o que dá autoridade para responder sem hesitar e sem
prometer demais.

\chapter{O centro de massa ponderado}
\label{cap:centro}

\section{A pergunta}

Onde está o ``ponto de equilíbrio geográfico'' de uma atividade num dado ano, e
para onde esse ponto se desloca ao longo do tempo?

\section{A fórmula}

Seja $i = 1, \dots, 166$ o índice das AMCs, $(x_i, y_i)$ o centroide da unidade
$i$ em coordenadas projetadas métricas (EPSG:5880), e $w_{it}$ o peso da
unidade $i$ no ano $t$. O centro médio ponderado é

\begin{equation}
\bar{x}_t = \frac{\sum_{i=1}^{n} w_{it}\, x_i}{\sum_{i=1}^{n} w_{it}},
\qquad
\bar{y}_t = \frac{\sum_{i=1}^{n} w_{it}\, y_i}{\sum_{i=1}^{n} w_{it}}.
\label{eq:centro}
\end{equation}

É, literalmente, a fórmula do centro de gravidade da física, com a massa
substituída pela quantidade de interesse. O deslocamento entre dois anos é o
vetor $(\bar{x}_T - \bar{x}_0,\ \bar{y}_T - \bar{y}_0)$, cuja componente $y$ é o
$\Delta$Norte reportado em quilômetros.

\section{Os dois tipos de peso, que respondem a perguntas diferentes}

A mesma fórmula recebe dois pesos, e confundi-los seria um erro grave de
leitura.

\textbf{Peso de estoque:} $w_{it}$ é a quantidade que a unidade \emph{guarda}
no ano --- a área da classe de cobertura, o efetivo do rebanho. Responde a
``onde a atividade está''. É de onde saem os números da tabela dos centros de
massa.

\textbf{Peso de fluxo:} $w_{it}$ é a área que, naquela unidade, \emph{mudou} de
uma classe para outra entre dois anos consecutivos, tomada da matriz de
transição e rotulada pelo ano de origem do par. Responde a ``onde a atividade
muda''. É de onde sai o centro da conversão.

\begin{armadilha}[A pergunta sobre originalidade]
A versão de \emph{estoque} é procedimento clássico --- é o centro de gravidade
de sempre, aplicado a outra variável, e nada tem de novo. É para a versão de
\emph{fluxo}, aplicada a matrizes de transição de uso da terra, que o
levantamento não localizou precedente. Se a banca perguntar o que há de
original no método, essa é a resposta honesta, e ela vem com a ressalva de que
a busca segue registrada como item em aberto. Enquanto estiver em aberto, a
métrica é defendida pela verificação interna, e não por filiação.
\end{armadilha}

\section{Os dois acompanhantes}

O \textbf{centro mediano} minimiza a soma das distâncias euclidianas aos
pontos, em vez da soma dos quadrados:
\begin{equation}
(\hat{x}, \hat{y}) = \arg\min_{(x,y)} \sum_{i} w_i \sqrt{(x - x_i)^2 + (y - y_i)^2}.
\end{equation}
Não tem solução fechada, e é obtido pelo algoritmo iterativo de Weiszfeld. A
razão de usá-lo é que ele é \textbf{robusto a um aglomerado dominante}: se uma
única região concentrasse massa desproporcional, ela puxaria o centro médio
mais do que puxa o mediano. Concordância entre os dois é evidência de que o
resultado não é artefato de concentração.

A \textbf{elipse de desvio-padrão} descreve dispersão e orientação da mancha, e
entra como procedimento de uso corrente.

\section{O que a decisão D19 impõe}

Um centro de massa é uma estatística pontual. Sem margem de erro, não há como
distinguir deslocamento real de ruído. Daí a regra:

\begin{quote}\itshape
Todo deslocamento vem com intervalo de confiança de 95\%, e \textbf{um
deslocamento cujo intervalo inclui zero nunca é reportado em quilômetros}. A
variável é descrita, nesse caso, como \emph{ancorada}.
\end{quote}

É por isso que a vegetação natural aparece como ``ancorada'' e não como ``andou
7,6 km''. Essa é uma autocorreção, e vale apresentá-la como tal.

\begin{ancora}[A tabela dos centros de massa]
Pastagem $+77{,}6$ km, IC $[+54{,}7;\ +98{,}2]$ --- robusto \\
Rebanho bovino $+66{,}9$ km, IC $[+47{,}2;\ +84{,}5]$ --- robusto \\
Agricultura $+65{,}2$ km, IC $[+43{,}5;\ +94{,}6]$ --- robusto \\
Vegetação natural $+7{,}6$ km, IC $[-0{,}5;\ +15{,}6]$ --- \textbf{ancorada}
\end{ancora}

\begin{naodiz}
O centro de massa não diz \emph{por que} a atividade se moveu, nem que uma
região tenha causado o movimento da outra. Ele é descrição. Toda a Parte da
explicação vem depois, e com outros instrumentos.
\end{naodiz}

\section{A régua de escala: o teste do MAUP}

O centro calculado sobre centroides de AMC supõe distribuição uniforme da
classe dentro de cada unidade --- suposição não neutra, porque as unidades do
Norte são maiores e mais irregulares que as do Sul. É o problema da unidade de
área modificável (MAUP): resultados calculados sobre unidades espaciais mudam
quando muda o recorte dessas unidades.

O teste foi direto: recalcular o mesmo centro \textbf{sem malha administrativa
nenhuma}, com cada pixel de 30 m pesando por si.

\begin{ancora}[Pixel contra AMC]
Pastagem: $+79{,}2$ km (pixel) contra $+77{,}6$ km (AMC). \\
Agricultura: $+66{,}9$ km (pixel) contra $+65{,}2$ km (AMC). \\
As duas medições concordam dentro de dois quilômetros, contra deslocamentos
medidos em dezenas.
\end{ancora}

\begin{nabanca}
``O MAUP existe no desenho, e eu o medi em vez de supor que fosse pequeno.
Refiz o cálculo sem malha nenhuma, pixel a pixel, e a diferença é de um a dois
quilômetros contra um deslocamento de setenta e oito. Não é material para esta
métrica.''
\end{nabanca}

\chapter{Bootstrap, e por que ele precisou virar bootstrap de blocos}
\label{cap:bootstrap}

\section{O bootstrap comum}

O problema: a fórmula \ref{eq:centro} devolve um número, não uma distribuição.
De onde sai o intervalo de confiança?

A resposta é reamostragem. O procedimento é o seguinte:

\begin{enumerate}
\item Sorteie, \emph{com reposição}, 166 AMCs do conjunto das 166. Algumas
aparecerão duas vezes, outras nenhuma.
\item Recalcule o centro de massa e o $\Delta$Norte sobre essa amostra.
\item Repita $B = 2000$ vezes.
\item O intervalo de 95\% são os percentis 2,5 e 97,5 da distribuição dos 2000
valores.
\end{enumerate}

A intuição é a seguinte: o bootstrap pergunta \emph{``se eu tivesse sorteado um
Goiás ligeiramente diferente, quanto o resultado mudaria?''}. Se ele muda pouco,
o número é estável; se muda muito, a estatística pontual estava escondendo
incerteza.

\section{Por que ele estava errado, e como foi corrigido}

O bootstrap comum supõe que as unidades são \textbf{trocáveis e
independentes}. O diagnóstico espacial do próprio trabalho mostra que não são:
o índice de Moran nos resíduos é significativo em 115 de 140 testes na malha
municipal e em 125 de 140 na malha de AMC.

Onde há dependência entre vizinhas, o número \emph{efetivo} de unidades
independentes é menor que 166, e o intervalo sai estreito demais. Sortear duas
AMCs vizinhas que dizem quase a mesma coisa não é o mesmo que sortear duas
informações.

A correção é o \textbf{bootstrap de blocos espaciais}: em vez de sortear AMC a
AMC, o estado é particionado em blocos espacialmente coerentes (agrupamento por
$k$-médias sobre os centroides, o que os torna compactos sem impor
contiguidade), e sorteiam-se os \emph{blocos inteiros}, com todos os seus
vizinhos juntos.

\section{O ponto sutil: o tamanho do bloco não é um valor a descobrir}

Este é o detalhe que separa uma resposta boa de uma ótima. O tamanho do bloco
não é um parâmetro a estimar, e sim um parâmetro que \emph{se varre}, porque
os dois extremos são conhecidos de antemão: um bloco de uma AMC reproduz
exatamente o bootstrap original, e blocos grandes dão intervalos cada vez mais
conservadores e menos precisos.

O que se reporta, portanto, não é um tamanho eleito, e sim a
\textbf{concordância do veredito} ao longo da grade de tamanhos.

\begin{ancora}[O resultado da varredura]
Tamanho de bloco varrido de \textbf{1 a 14} AMCs. No caso da pastagem, o
intervalo alarga de \textbf{45 para 80 km} de largura, como esperado. \\
\textbf{O veredito não muda em nenhum tamanho:} pastagem, rebanho e agricultura
excluem zero nas seis partições testadas; a vegetação natural inclui zero nas
seis. \\
Ponto sensível: \textbf{nenhuma} componente \emph{leste} atravessa a varredura
inteira --- a da pastagem inclui zero já em blocos de três AMCs, a da
agricultura em oito, a do rebanho em catorze.
\end{ancora}

\begin{nabanca}
``Meu primeiro intervalo tratava as AMCs como independentes, e o meu próprio
diagnóstico espacial mostra que elas não são. Refiz com bootstrap de blocos
espaciais e o intervalo quase dobrou de largura, como tinha de acontecer. O
veredito não muda em nenhum tamanho de bloco, então a leitura da tabela não
depende de tratar as unidades como independentes.''
\end{nabanca}

\begin{naodiz}
O bootstrap não corrige viés do estimador nem erro do classificador. Ele mede
apenas a incerteza de amostragem sobre as unidades. A acurácia do MapBiomas
está fora do alcance de qualquer reamostragem, e isso é declarado.
\end{naodiz}

\chapter{Quebras estruturais: de onde vêm os três atos}
\label{cap:quebras}

\section{A pergunta, e por que ela é metodológica antes de ser narrativa}

Narrar quarenta anos exige dividi-los, e onde cai cada divisa muda o que a
narrativa mostra. Ancorar as divisas em marcos institucionais escolhidos de
antemão deixaria sem resposta a pergunta ``por que 2005 e não 2006?''. A lógica
adotada é a inversa: as fronteiras saem dos dados, e os marcos institucionais
são classificados \emph{depois}, pela força da evidência que os sustenta.

\section{O método primário: a estatística sup-F de Quandt-Andrews}

A ideia é simples e vale ser dita em português antes da fórmula. Suponha que
você queira testar se houve uma quebra no ano $\tau$. Você estima dois modelos
--- um só, para a série inteira, e dois, um antes e um depois de $\tau$ --- e
compara o quanto o segundo arranjo melhora o ajuste. Isso é um teste F comum.

O problema é que você não sabe $\tau$ de antemão. A solução de Quandt e Andrews
é: calcule o F para \emph{todos} os anos-candidatos e fique com o maior.

\begin{equation}
\text{sup-}F = \max_{\tau \in [\tau_{\min},\, \tau_{\max}]} F(\tau).
\end{equation}

Aqui a estatística é \textbf{multivariada}, isto é, aplicada conjuntamente às
variações anuais de vegetação natural, pastagem e agricultura, com segmentação
binária.

\subsection{Os três parâmetros que precisam ser declarados}

Eles determinam quais anos-candidatos o método chega a devolver, e foram
fixados \emph{antes} de rodar os testes:

\begin{itemize}
\item \textbf{Segmento mínimo de cinco anos} --- o \emph{trimming} que o sup-F
exige. Sem ele, a estatística explode nas bordas da série, onde um dos dois
segmentos tem pouquíssimas observações.
\item \textbf{Máximo de três quebras} procuradas.
\item \textbf{Limiar de $F \geq 4$} para aceitar uma quebra nova --- mais
frouxo que o usual no caso univariado, porque a estatística aqui é conjunta a
três séries.
\end{itemize}

\section{Os dois métodos de sensibilidade}

\textbf{O detector sequencial de Rodionov (STARS)} percorre a série testando,
a cada ponto, se a média mudou o suficiente para declarar um novo regime.
Entra aqui porque ele alcança regimes \emph{curtos} que o \emph{trimming} de
cinco anos do sup-F descarta por construção.

\begin{armadilha}[A honestidade que vale ouro na banca]
A implementação usada \textbf{não reproduz o algoritmo publicado em todos os
pontos}: mantém o teste sequencial de médias, mas dispensa a etapa de
confirmação pelo índice de mudança de regime, estima a dispersão de forma
robusta sobre a série inteira e adota limiar mais exigente que o do artigo. As
três escolhas puxam em direções que não se somam de modo previsível, e o efeito
líquido não foi medido. \textbf{É por isso que o detector entra apenas como
sensibilidade}, com poder de corroborar ou de deixar de corroborar uma
fronteira que o primário já apontou, nunca de estabelecer fronteira por conta
própria. Dizer isso antes de a banca perguntar transforma uma vulnerabilidade
em demonstração de rigor.
\end{armadilha}

\textbf{A divergência entre matrizes de transição} compara a matriz média
anterior a cada ano-candidato com a posterior a ele. Duas medidas, de alcance
desigual:

\begin{itemize}
\item a divergência de \textbf{Kullback-Leibler} entre as duas matrizes,
ponderada pela distribuição estacionária, com significância por permutação. Ela
cresce ao longo de quase toda a série, e por isso informa mais sobre
\emph{direção} do que sobre a \emph{localização} de uma fronteira;
\item a distância de \textbf{variação total} entre as distribuições
estacionárias que essas matrizes implicam, que localiza melhor, mas não recebe
teste de significância e entra como leitura auxiliar.
\end{itemize}

Para duas distribuições $P$ e $Q$ sobre as mesmas classes:
\begin{equation}
D_{\mathrm{KL}}(P \| Q) = \sum_{k} P_k \log \frac{P_k}{Q_k},
\qquad
d_{\mathrm{TV}}(P, Q) = \tfrac{1}{2}\sum_{k} |P_k - Q_k|.
\end{equation}

\section{A regra de decisão, fixada antes dos testes}

Uma fronteira só é confirmada quando o método primário \emph{e ao menos um}
método de sensibilidade convergem. Divergência entre os três seria documentada
como sensibilidade ao método, nunca resolvida por escolha.

\begin{ancora}[O que o procedimento devolveu]
\textbf{2001} --- $F = 62{,}2$ (adotada) \\
\textbf{2020} --- $F = 21{,}5$ (adotada) \\
\textbf{1991} --- $F = 10{,}3$ (suporte bem mais fraco; \textbf{não adotada}) \\[4pt]
Daí os três atos: \textbf{1985--2000}, \textbf{2001--2019}, \textbf{2020--2024}.
\end{ancora}

\section{As duas validações independentes, que são o ponto forte}

\textbf{Primeira: a âncora externa.} O p-valor do sup-F é bruto, porque muitos
anos-candidatos são testados. A fronteira de 2020 é sustentada, decisivamente,
por uma fonte que \emph{não passa por classificador algum}: a área de soja
plantada que o IBGE recolhe em campo quebra em 2020 sozinha.

\textbf{Segunda: a verificação de falso positivo.} O mesmo procedimento foi
aplicado a séries de \emph{ruído}, para medir quantas quebras aparecem por
acaso.

\begin{naodiz}
O teste de Kruskal-Wallis que mostra que os três atos diferem em taxa de mudança
($H = 20{,}3$; $p < 0{,}001$) \textbf{descreve} o contraste; ele não o valida,
porque as fronteiras foram estimadas dos mesmos dados. Essa circularidade é
declarada no texto, e a validação vem de fora.
\end{naodiz}

\section{A ausência de quebra também é resultado}

O Código Florestal de 2012 não produz quebra empírica em nenhuma das três
séries. Isso rebaixa o seu estatuto de marco \emph{periodizador}, e não diz
nada sobre seus efeitos por outras vias. A distinção é fina e vale ser dita com
essas palavras.

\section{O quarto ato que não foi adotado}

Uma candidata em torno de 2005--2006 foi examinada e recusada, por três razões
que vale ter prontas:

\begin{enumerate}
\item o método primário não a detecta, porque o sup-F multivariado lê o
intervalo entre as duas fronteiras confirmadas como platô contínuo, e a regra
exige convergência;
\item a taxa total de mudança não separa as duas sub-fases: comparadas
2001--2005 e 2006--2019, o teste devolve $p = 0{,}060$, com poder de
\textbf{0,63} para as quatro observações disponíveis;
\item a fronteira é sensível ao ponto de corte: significativa em 2005
($p = 0{,}046$), marginal em 2004 ($p = 0{,}10$) e nula em 2006 ($p = 0{,}12$).
\end{enumerate}

O que a sub-fase tem de real é \emph{composicional}, e fica registrado como nota
dentro do Ato~II: nela a vegetação natural se perde em ritmo bem mais alto do
que em 2006--2019 ($p = 0{,}0008$), e a conversão de pastagem em agricultura é
cerca de quatro vezes mais intensa. É diferença de composição, não de
intensidade total.

\section{A convenção de contagem entre atos, e o teste que a acompanha}

Sempre que uma grandeza é medida \emph{entre} dois anos --- fluxos, emissão,
taxas por ato ---, cada ato é medido entre o seu primeiro e o seu último ano.
As duas transições de virada, 2000--2001 e 2019--2020, não pertencem a ato
nenhum. A alternativa faria a conta do ato depender de uma escolha sem
critério, e justamente no ponto em que o regime muda.

O custo é que a soma dos atos fica abaixo do período inteiro --- 3\% na conta de
carbono. E o efeito dessa escolha \textbf{foi medido, e não suposto}:

\begin{ancora}[A sensibilidade da convenção]
Parcela do Ato~I no estoque removido: \textbf{76,4\%} na convenção adotada,
76,6\% se a virada vai para o ato anterior, 74,1\% se vai para o seguinte. \\
Parcela do Ato~II: 18,6\%, 18,6\% e 20,5\%. \\
A repartição se move até 2,3 pontos, e a afirmação de ``três quartos no Ato~I''
sobrevive às três convenções.
\end{ancora}

\chapter{Matrizes de transição e o que é uma cadeia de Markov aqui}
\label{cap:markov}

\section{A construção}

Para um par de anos $(t, t+1)$, a matriz de transição $M$ tem entradas
$M_{jk} = $ número de pixels que estavam na classe $j$ em $t$ e estão na classe
$k$ em $t+1$. Normalizada por linha, ela vira uma matriz de probabilidades de
transição:
\begin{equation}
P_{jk} = \frac{M_{jk}}{\sum_{k'} M_{jk'}},
\qquad \sum_k P_{jk} = 1.
\end{equation}

A matriz primária deste trabalho opera sobre \textbf{sete grupos de classes},
incluindo o Mosaico de Usos como grupo próprio.

\section{A distribuição estacionária, e para que ela serve aqui}

Se as probabilidades $P$ permanecessem constantes indefinidamente, a paisagem
convergiria para uma distribuição $\pi$ que satisfaz
\begin{equation}
\pi = \pi P, \qquad \sum_k \pi_k = 1,
\end{equation}
isto é, o autovetor à esquerda de $P$ associado ao autovalor 1.

\begin{naodiz}
Esta $\pi$ \textbf{não é uma previsão} do que Goiás será. Ela é o resumo do
regime que aquela matriz descreve --- ``para onde esta dinâmica levaria se ela
não mudasse''. Como a dinâmica muda (é justamente o que a periodização detecta),
$\pi$ entra aqui apenas como \emph{descritor comparável entre períodos}, para
alimentar a distância de variação total do capítulo anterior. Nenhuma projeção
do trabalho se apoia nela.
\end{naodiz}

\section{Por que a matriz não pode ser inferida da diferença entre estoques}

Ponto conceitual que vale dizer devagar. Se a pastagem perdeu 1 Mha e a lavoura
ganhou 1 Mha, é tentador concluir que 1 Mha de pasto virou lavoura. Mas o mesmo
saldo é compatível com 3 Mha de pasto virando lavoura enquanto 2 Mha de Cerrado
viram pasto. Só a contagem pixel a pixel decompõe cada hectare que mudou em um
fluxo \emph{orientado}, e é ela que autoriza a afirmação de origem.

\section{O que o cruzamento entre as duas pontas não vê}

A figura que cruza 1985 com 2024 registra o estado \emph{inicial} e o
\emph{final} de cada pixel, e não as passagens intermediárias. Um pixel que foi
pasto, virou lavoura e voltou a pasto entra como pasto. Por isso os fluxos do
cruzamento de pontas são \emph{menores} que a contagem de eventos anuais.

\begin{ancora}[Os quatro fluxos que desenham a marcha, 1985$\to$2024]
Vegetação natural $\to$ pastagem: \textbf{4,10 Mha} (o maior da série) \\
Pastagem $\to$ agricultura: \textbf{2,72 Mha} (a maior entre classes já manejadas) \\
Vegetação natural $\to$ agricultura: \textbf{1,29 Mha} \\
Vegetação natural $\to$ Mosaico de Usos: \textbf{1,00 Mha} (mantido \emph{fora}
da conta de conversão antrópica) \\[4pt]
A pastagem é o elo intermediário de \textbf{58\%} da entrada em agricultura
(2,72 de 4,65 Mha). Entre as duas rotas vindas de fora do sistema já manejado,
a proporção sobe a \textbf{68\%}.
\end{ancora}

\begin{nabanca}[Por que o quarto fluxo ficou de fora]
``Incluir a conversão de vegetação natural em Mosaico faria a minha medida de
desmatamento saltar 35\% acima do PRODES. E como essa classe tem transição de
ida e de volta, tratá-la como perda definitiva superestima. Mantê-la fora é a
escolha conservadora, e ela está declarada.''
\end{nabanca}

\section{As duas convenções de classe que precisam ser sabidas de cor}

Há dois pontos em que a tabela do balanço e a figura do cruzamento usam
agrupamentos diferentes. A diferença é de \emph{agrupamento}, não de dado, e a
banca pode cobrar a reconciliação:

\begin{itemize}
\item \textbf{Vegetação natural.} A tabela do balanço soma formação florestal,
savânica, campo nativo \emph{e campo alagado}, e dá 17,65 Mha em 1985. O
agrupamento em sete classes do censo de pixels põe o campo alagado em
``Outros'', e dá 16,96 Mha. A diferença é exatamente o campo alagado,
\textbf{0,70 Mha}.
\item \textbf{Agricultura.} Na tabela do balanço ela é a lavoura \emph{sem}
silvicultura, e dá 5,58 Mha em 2024. No agrupamento em sete classes ela a
inclui, e dá 5,73 Mha. A diferença é a silvicultura, \textbf{0,16 Mha}.
\end{itemize}

\begin{nabanca}
``Não reagrupei o cubo de propósito, porque mover a classe do campo alagado
mexeria nos fluxos-manchete de 4,10, 2,72 e 1,29 Mha, que são auditados e
citados. Em vez disso, declarei a composição em nota de rodapé e as duas
legendas apontam uma para a outra.''
\end{nabanca}

\chapter{Mistura de gaussianas: duas populações de pastagem}
\label{cap:mistura}

\section{A pergunta}

Quando um pixel de pastagem é convertido em lavoura, quantos anos ele tinha de
pastagem? E a distribuição dessas idades é composta por um só tipo de pastagem
ou por mais de um?

A pergunta importa porque a idade da pastagem no instante da conversão é a
aproximação mais próxima da \emph{intenção do produtor} que a série permite
observar. Um pasto de três anos que vira lavoura é etapa planejada de um
sistema; um pasto de trinta anos que vira lavoura é reserva de terra sendo
ativada. São duas economias diferentes com o mesmo nome no mapa.

\section{A fórmula}

O modelo de mistura de duas gaussianas supõe que a densidade das idades $a$ é
\begin{equation}
f(a) = \pi_1 \, \mathcal{N}(a \mid \mu_1, \sigma_1^2)
     + \pi_2 \, \mathcal{N}(a \mid \mu_2, \sigma_2^2),
\qquad \pi_1 + \pi_2 = 1,
\end{equation}
com cinco parâmetros livres ($\mu_1, \sigma_1, \mu_2, \sigma_2, \pi_1$),
estimados por máxima verossimilhança via algoritmo EM. Contra ele se ajusta o
modelo de uma componente só, com dois parâmetros.

\begin{ancora}[O resultado]
População jovem: $\mu \approx 4$ anos, $\sigma \approx 1{,}6$ ano,
\textbf{33\%} da massa. \\
População velha: $\mu \approx 16$ anos, $\sigma \approx 7{,}5$ anos,
\textbf{67\%} da massa, com cauda que alcança os quarenta anos da série. \\
Uma componente sozinha devolve $\mu \approx 12$ anos e \textbf{erra as duas
pontas}, subestimando o pico jovem e a cauda velha.
\end{ancora}

\section{Por que a curva não tem dois picos, e por que isso não é problema}

Detalhe que a banca pode levantar olhando a figura. A forma não tem dois picos
visíveis, e sim um pico jovem e um \textbf{ombro} longo. A razão é aritmética:
sendo cerca de cinco vezes mais dispersa que a primeira ($7{,}5$ contra
$1{,}6$), a segunda população espalha a mesma massa por uma faixa muito maior e
produz platô, e não segunda espiga.

\begin{nabanca}
``A separação não se lê por inspeção visual, e eu não a estabeleço assim. O que
a estabelece é que uma população única não dá conta da forma: ajustada sozinha
ela erra as duas pontas.''
\end{nabanca}

\section{O problema do censo, e por que o BIC foi rebaixado (decisão D23)}

Este é um dos pontos mais sofisticados do trabalho, e um que impressiona quando
bem explicado.

O critério de informação bayesiano é
\begin{equation}
\mathrm{BIC} = -2 \ln \hat{L} + k \ln n,
\end{equation}
onde $\hat{L}$ é a verossimilhança maximizada, $k$ o número de parâmetros e $n$
o número de observações. Compara-se o BIC dos dois modelos, e o menor vence.

O problema é o seguinte. A análise roda sobre o \textbf{censo} de eventos, e não
sobre uma amostra: $n$ é a população inteira, e são \textbf{16,0 milhões} de
eventos com idade conhecida. Com $n$ dessa ordem, qualquer afastamento ínfimo do
modelo de uma componente, e nenhum dado real é exatamente gaussiano ---
produz diferença de BIC astronômica.

\begin{quote}\itshape\color{cortinta}
A magnitude do $\Delta$BIC mede o tamanho de $n$, e não a força da evidência.
\end{quote}

Por isso o $\Delta$BIC \textbf{não} é usado como critério probatório aqui. E o
que sustenta a leitura de duas populações são evidências que \emph{não dependem
de ajuste algum}.

\section{A evidência que não depende do modelo}

Separados apenas pela origem anterior do pixel, os eventos se ordenam sozinhos:

\begin{ancora}[As duas pontas, sem ajustar nada]
Pasto que veio de \textbf{lavoura} volta a ser lavoura com mediana de
\textbf{5 anos}. \\
Pasto que veio de \textbf{vegetação natural} dura o dobro: mediana de
\textbf{13 anos}. \\[4pt]
Nas conversões de pasto com até \textbf{três anos}, \textbf{45\%} já tinham sido
lavoura antes --- lavoura, pasto, lavoura: rotação. \\
Nas de pasto com \textbf{trinta anos ou mais}, a lavoura anterior praticamente
desaparece: \textbf{1\%} --- Cerrado, pasto, lavoura: reserva ativada.
\end{ancora}

O contraste de 45\% contra 1\% é o que o dado estabelece com firmeza. A
atribuição de cada pixel individual, não.

\section{O critério triplo de separação}

Onde o ajuste é usado --- para verificar se as duas populações aparecem em cada
recorte ---, o critério é declarado e é triplo:

\begin{enumerate}
\item o ajuste de duas componentes é preferido pelo critério de informação;
\item os modos estão afastados em \textbf{pelo menos cinco anos};
\item a componente menor responde por \textbf{ao menos 15\%} da massa.
\end{enumerate}

Sob censo, quem decide são as \emph{duas últimas} exigências, pela razão acima.
E nenhuma delas é a presença de um vale na densidade, que a curva agregada não
tem.

\begin{naodiz}
A mistura não prova que a população jovem seja integração lavoura-pecuária. Ela
é \emph{compatível} com isso e não a comprova, porque o satélite vê o capim e
não vê se havia gado nele. \\[4pt]
E o Mosaico de Usos é a classe em que o classificador não separou usos: metade
da origem do pasto velho fica, por construção, indeterminada entre vegetação
natural e uso misto anterior.
\end{naodiz}

\begin{naodiz}[E nada se afirma quanto a tendência]
O eixo do tempo está comprometido nas duas pontas. Antes, pela idade truncada no
início da série; depois, pela mudança de rótulo. Que o Ato~I apareça com uma
população só \textbf{não} significa que a reserva antiga tenha surgido depois:
uma conversão em 1995 ocorre sobre um pasto que tem, no máximo, dez anos, porque
a série começa em 1985. Para observar a conversão de um pasto de dezesseis anos
--- o modo da população velha --- é preciso chegar a 2001. No primeiro período,
a população velha é \textbf{invisível, não ausente}.
\end{naodiz}

\section{As duas objeções testadas, e por que o achado saiu mais forte}

A curva soma quarenta anos de conversões em todo o estado. Talvez não haja duas
populações convivendo em lugar algum, e sim regiões ou épocas diferentes
somadas. As duas objeções foram testadas isolando cada recorte.

\textbf{Não são regiões somadas.} As cinco mesorregiões separam-se em duas
populações, e também o fazem todas as 166 AMCs na régua imune à mudança de
rótulo (na régua estrita, 162 das 164 que registram conversão suficiente, e as
duas exceções falham pelo mesmo lado). A mesorregião explica \textbf{0,5\%} da
variação da idade.

\textbf{Não são épocas somadas.} Isoladas por período, as duas populações
permanecem: no Ato~II sozinho (32,5 milhões de conversões pela régua da união, o
maior bloco da série) e no Ato~III sozinho, nas cinco mesorregiões de cada um,
em \textbf{10 de 10} células de período por região (9 de 10 pela régua estrita).

\begin{armadilha}[Dois números que parecem contraditórios e não são]
Os 32,5 milhões do Ato~II \textbf{não} se somam aos 16,0 milhões de eventos com
idade conhecida. Aqueles são eventos na régua da união; estes são eventos com
idade conhecida na régua estrita. As duas contagens medem populações
diferentes. Se a banca cruzar os números, essa é a resposta.
\end{armadilha}

\begin{nabanca}[Como apresentar a autocorreção mais extensa do trabalho]
``A versão que eu teria publicado dizia que o Sul convertia pasto jovem, com
mediana de nove anos, e o Norte pasto velho, com dezesseis. O defeito não estava
na conta, e sim no que entrava nela: a idade só é medida quando o destino do
pixel recebe o rótulo de agricultura, e no fim da série parte da conversão passa
a ser rotulada como Mosaico. Esses eventos saíam da amostra sem aviso, e não
saíam por igual no território --- concentravam-se mais ao norte. O gradiente
era, em boa parte, essa seleção. Refeita a análise com a pergunta grossa, a
amplitude Sul--Norte cai de sete anos para dois, a ordem das regiões se
embaralha e a parcela de variação atribuível à mesorregião despenca de 3,7\%
para 0,5\%. E o que sobrevive é mais forte do que o que caiu: quanto menos a
região explica, mais universal se mostra a convivência dos dois mecanismos.''
\end{nabanca}
